\documentclass[12pt,a4paper]{article}
\usepackage{ugmath}
\usepackage{float}
\usepackage{placeins}
\usepackage [spanish]{babel}
\usepackage[labelfont=bf,labelsep=space]{caption}
\newcommand{\paren}[1]{\left( #1 \right)}
\newcommand{\alumno}{Ricardo León Martínez}
\newcommand{\materia}{Calculo Diferencial e Integral II}
\newcommand{\profesor}{Fernando Nuñez Medina}
\newcommand{\tarea}{Tarea 9}
\newcommand{\fecha}{25/4/2026}

\begin{document}
    \begin{center}
    {\large\textbf{UNIVERSIDAD DE GUANAJUATO}}\\[0.3cm]
    {\normalsize\textbf{DIVISIÓN DE CIENCIAS NATURALES Y EXACTAS}}\\
    {\normalsize\textbf{CAMPUS GUANAJUATO}}\\[1cm]

    {\Large\textbf{\tarea\ (\materia)}}\\[1cm]
    \end{center}

    \textbf{Nombre:} \alumno \hfill 
    \textbf{Fecha:} \fecha \hfill 
    \textbf{Calificación:} \rule{3cm}{0.4pt} \\[0.3cm]

    \begin{mdframed}[style=mdpinkbox,frametitle={Ejercicio 1}]
        Prueba el corolario 5.
    \end{mdframed}

    \begin{mdframed}[style=mdpinkbox,frametitle={Ejercicio 2}]
        \textbf{(Limites y desigualdades)} Considera dos sucesiones convergentes $(a_{n})$  y $(b_{n})$.
        \begin{enumerate}
            \item[(a)] Prueba que si existe $N\in\N$ tal que
            \begin{equation*}
                a_{n}\leq b_{n}\text{ para todo } n\gt N
            \end{equation*}
            entonces
            \begin{equation*}
                \lim_{n\to\infty}a_{n}\leq\lim_{n\to\infty}b_{n}.
            \end{equation*}
            \item[(b)] ¿Se cumplira que si $a_{n}\lt b_{n}$ para todo $n\gt N$, entonces $\lim_{n\to\infty}a_{n}\lt\lim_{n\to\infty}b_{n}$? Argumenta tu respuesta. 
        \end{enumerate}
    \end{mdframed}

    \begin{proof}
        \begin{enumerate}
            \item[(a)]
            Procedemos por contradicción. Supongamos que $a\gt b$. Definimos
            \begin{equation*}
                \varepsilon:=\frac{a-b}{2}\gt0.
            \end{equation*}
            Por la convergencia de las sucesiónes existen $N_{1},N_{2}\in\N$ tales que
            \begin{equation*}
                \text{si }n\gt N_{1},\text{ entonces }\abs{a_{n}-a}\lt\varepsilon,\qquad \text{si }n\gt N_{2},\text{ entonces }\abs{b_{n}-b}\lt\varepsilon.
            \end{equation*}
            Sea $N:=\max\{N_{0},N_{1},N_{2}\}$. Entonces, para todo $n\gt N$, se cumple
            \begin{equation*}
                a_{n}\gt a-\varepsilon=\frac{a+b}{2},\qquad b_{n}\lt b+\varepsilon=\frac{a+b}{2}.
            \end{equation*}
            Por lo tanto,
            \begin{equation*}
                a_{n}\gt b_{n},
            \end{equation*}
            lo cual contradice la hipotesis $a_{n}\leq b_{n}$ para todo $n\gt N_{0}$.
            \item[(b)]
            Esto es falso en general. Consideremos
            \begin{equation*}
                a_{n}=0,\qquad b_{n}=\frac{1}{n}.
            \end{equation*}
            Entonces $a_{n}\lt b_{n}$ para todo $n$, pero
            \begin{equation*}
                \lim_{n\to\infty}a_{n}=0,\qquad \lim_{n\to\infty}b_{n}=0.
            \end{equation*}
            Por lo tanto,
            \begin{equation*}
                \lim_{n\to\infty}a_{n}=\lim_{n\to\infty}b_{n}.
            \end{equation*}
        \end{enumerate}
    \end{proof}

    \begin{mdframed}[style=mdpinkbox,frametitle={Ejercicio 3}]
        Prueba que si una sucesion es de Cauchy posee una subsucesion convergente entonces converge.
    \end{mdframed}

    \begin{proof}
        Veamos que la sucesión es de cauchy, por la proposición 54 esta converge.
    \end{proof}

    \begin{mdframed}[style=mdpinkbox,frametitle={Ejercicio 4}]
        Prueba que toda sucesion de numeros reales tiene un limite subsecuencial en $\R^{*}$.
    \end{mdframed}

    \begin{proof}
        Sea $(a_{n})$ una sucesión de números reales. Consideremos tres casos. Primero $(a_{n})$ es acotada.
        Entonces por el teorema de Bolzano-Weierstrass, existe una subsucesión $(a_{n_{k}})$ convergente en $\R$,
        es decir,
        \begin{equation*}
            a_{n_{k}}\to\ell\in\R.
        \end{equation*}
        Como $\R\subset\R^{*}$, esta subsucesión tmb converge en $\R^{*}$. Segundo $(a_{n})$ no esta acotada superiormente.
        Entonces, para todo $M\in\R$, existe $n$ tal que $a_{n}\gt M$. Construimos una subsecesión por inducción. Elegimos
        $n_{1}$ tal que $a_{n_{1}}\gt1$. Elegimos $n_{2}\gt n_{1}$ tal que $a_{n_{2}}\gt2$. En general, elegimos $n_{k}\gt n_{n_{k-1}}$
        tal que $a_{n_{k}}\gt k$. Esto es posible porque la sucesión no está acotada superiormente. Entonces
        \begin{equation*}
            a_{n_{k}}\gt k,\quad \forall k\in\N.
        \end{equation*}
        Por lo tanto,
        \begin{equation*}
            a_{n_{k}}\to+\infty,
        \end{equation*}
        es decir, converge en $+\infty$. Tercero $(a_{n})$ no esta acotada inferiormente. Analogamente, para todo $M\in\R$,
        existe $n$ tal que $a_{n}\lt M$. Volvemos a construir una subsucesión por inducción. Elegimos $n_{1}$ tal que $a_{n_{1}}\lt-1$,
        $n_{2}\gt n_{1}$ tal que $a_{n_{2}}\lt-2$, en general $n_{k}\gt n_{k-1}$ tal que $a_{n_{k}}\lt -k$. Entonces
        \begin{equation*}
            a_{n_{k}}\lt-k,\quad\forall k,
        \end{equation*}
        y por lo tanto
        \begin{equation*}
            a_{n_{k}}\to-\infty.
        \end{equation*}
        En todos los casos, existe una subsucesión que converge en $\R^{*}$.
    \end{proof}

    \begin{mdframed}[style=mdpinkbox,frametitle={Ejercicio 5}]
        Prueba que si $(a_{n})$ es una sucesión de números reales, $N$ es un número natural y $(b_{n})$ es la suecesión
        $(a_{N},a_{N+1},a_{N+2},\dots)$, entonces $\overline{lim}b_{n}=\overline{lim}a_{n}$. Se cumple un resultado similar
        para el limite inferior.
    \end{mdframed}

    \begin{mdframed}[style=mdpinkbox,frametitle={Ejercicio 6}]
        Prueba que si $(a_{k})$ es una subsucesión de una sucesión $(a_{n})$ de números reales, entonces
        \begin{equation*}
            \liminf_{n\to\infty}a_{n}\leq\liminf_{j\to\infty}a_{n_{k}}\leq\limsup_{k\to\infty}a_{n_{k}}\leq\limsup_{n\to\infty}a_{n}.
        \end{equation*}
    \end{mdframed}

    \begin{proof}
        Sea $(a_{n})$ una sucesión en $\R$, y $(a_{n_{k}})$ una subsucesión, con $n_{k}$ estrictamente creciente. Definimos
        \begin{equation*}
            s_{n}:=\sup\{a_{m}:m\geq n\},\qquad i_{n}=\{a_{m}:m\geq n\}.
        \end{equation*}
    \end{proof}

    \begin{mdframed}[style=mdpinkbox,frametitle={Ejercicio 7}]
        Sea $m$ un número natural. Prueba que $\sum_{n=1}^{\infty}a_{n}$ converge ssi $\sum_{n=m}^{\infty}a_{n}$ converge;
        y que si $m\gt 1$,
        \begin{equation*}
            \sum_{n=1}^{\infty}a_{n}=(a_{1}+a_{2}+\cdots+a_{m-1})+\sum_{n=m}^{\infty}a_{n}.
        \end{equation*}
        Lamaremos al número $C_{m}:=\sum_{n=m}^{\infty}a_{n}$ la $m$-ésima cola de la serie $\sum_{n=1}^{\infty}a_{n}$.
        Así, una series converge ssi todas sus colas convergen.
    \end{mdframed}

    \begin{mdframed}[style=mdpinkbox,frametitle={Ejercicio 8}]
        Prueba que si  una serie $\sum_{n=1}^{\infty}a_{n}$ converge, entonces
        \begin{equation*}
            \lim_{m\to\infty}\sum_{n=m}^{\infty}a_{n}=0.
        \end{equation*}
        Así, una serie converge, entonces la sucesión formada por sus colas converge a cero.
    \end{mdframed}

\end{document}